%%%%%%%%%%%%%%%%%%%%%%%%%%%%%%%%%%%%%%%%%%%%%%%%%%%%%%%%%%%%%%%%%%%%%%%%%%%
%
% Generic template for TFC/TFM/TFG/Tesis
%
% $Id: desarrollo.tex,v 1.4 2015/06/05 00:05:18 macias Exp $
%
% By:
%  + Javier Macías-Guarasa.
%    Departamento de Electrónica
%    Universidad de Alcalá
%  + Roberto Barra-Chicote.
%    Departamento de Ingeniería Electrónica
%    Universidad Politécnica de Madrid
% 
% Based on original sources by Roberto Barra, Manuel Ocaña, Jesús Nuevo, Pedro Revenga, Fernando Herránz and Noelia Hernández. Thanks a lot to all of them, and to the many anonymous contributors found (thanks to google) that provided help in setting all this up.
%
% See also the additionalContributors.txt file to check the name of additional contributors to this work.
%
% If you think you can add pieces of relevant/useful examples, improvements, please contact us at (macias@depeca.uah.es)
%
% You can freely use this template and please contribute with comments or suggestions!!!
%
%%%%%%%%%%%%%%%%%%%%%%%%%%%%%%%%%%%%%%%%%%%%%%%%%%%%%%%%%%%%%%%%%%%%%%%%%%%

\chapter{El Problema del Viajante}
\label{cha:desarrollo}

Antes de comenzar con el desarrollo de las diferentes metaheurísticas propuestas, es importante realizar una explicación 
previa sobre el problema en sí. El \ac{TSP} es uno de los problemas computacionales más estudiados de las últimas décadas.
Se puede formular de la siguiente manera: dada una lista de ciudades, y las distancias entre todas ellas, se debe buscar la 
ruta más corta que visite cada una de las ciudades y vuelva finalmente a la ciudad de inicio. Formalmente, se puede definir
como un grafo ponderado completo:

\begin{equation}
  \label{eq:eq1}
  G = (V, E, c)
\end{equation}

Donde $V$ representa el conjunto de vértices o nodos del grafo, y $E$ el conjunto de aristas que conectan cada par de vértices.
A cada arista $(i, j) \in E$ se le asocia un coste $c_{ij}$, que representa la distancia desde el vértice $i$ hasta el vértice 
$j$.

Sea $\pi = ({\pi}_{1}, {\pi}_{2}, \dots , {\pi}_{n})$ una permutación de vértices de $V$, el objetivo consiste en encontrar la 
permutación que minimiza el coste total del recorrido, sujeta a que cada vértice sea visitado una única vez antes de volver al 
punto de partida:

\begin{equation}
  \label{eq:eq2}
    \min_{\pi}
    \left\{
    \sum_{i=1}^{n-1} c_{\pi_i,\pi_{i+1}}
    +
    c_{\pi_n,\pi_1}
    \right\},\
    \text{donde} \ \pi_i \neq \pi_j,\ \forall\, i\neq j.
\end{equation}

En otras palabras, se pretende obtener el ciclo hamiltoniano de coste mínimo del grafo. Un ciclo hamiltoniano es un recorrido
cerrado que visita todos los vértices exactamente una vez y finaliza en el vértice en el que comenzó \cite{zambito2006}. La
figura~\ref{fig:ham_cyc} muestra un ejemplo sencillo de un ciclo hamiltoniano:

\begin{figure}[H]
  \centering 
  \includegraphics[width=0.3\linewidth]{hamiltonian-cycle.png}
  \caption{Ciclo hamiltoniano sobre un conjunto de cinco ciudades.}
  \label{fig:ham_cyc}
\end{figure}

\section{Variantes del problema}
\label{sec:variantes}

Es conveniente destacar que el TSP admite distintas variantes, en función de las características del grafo y de las 
restricciones al recorrido. Aunque todas ellas tienen el mismo objetivo, las diferencias entre estas formulaciones 
condicionan la dificultad del problema y las técnicas empleadas para resolverlo. 

La formulación más habitual es la del \ac{STSP}, donde el coste de desplazamiento entre dos ciudades es el mismo en ambos
sentidos, cumpliéndose:

\begin{equation}
  \label{eq:eq3}
    c_{ij} = c_{ji},\
    \forall\, (i,j)\in E.
\end{equation}

Esta situación es la habitual cuando se toman las distancias euclídeas entre las ciudades. Por el contrario, en el \ac{ATSP}
los costes dependen del sentido del desplazamiento, de modo que:

\begin{equation}
  \label{eq:eq4}
    c_{ij} \neq c_{ji},\
    \exists\,(i,j)\in E.
\end{equation}

Esta variante aparece con frecuencia en aplicaciones reales, por ejemplo: calles con un único sentido, diferencias de pendiente, 
o condiciones de tráfico.

Además de estas dos variantes, la literatura recoge otras diferentes para contextos específicos. Entre ellas destacan el 
\ac{TSPTW}, en el que cada ciudad debe visitarse dentro de un determinado intervalo de tiempo \cite{dumas1995}; o el \ac{MTSP},
donde varios viajantes colaboran para visitar el conjunto de ciudades \cite{bektas2006}. Otra variante más reciente pero muy
interesante es la del \ac{DG-TSP}, cuyo objetivo consiste en minimizar simultáneamente el coste del recorrido y el impacto 
medioambiental, considerando factores dinámicos como el tráfico, el gradiente de la carretera o el consumo energético 
del vehículo \cite{benjelloun2025}.

En este trabajo se va a considerar exclusivamente el Problema del Viajante simétrico, ya que se corresponde con
la formulación más clásica del problema y la más utilizada para evaluar y comparar algoritmos de optimización.

\section{Espacio de búsqueda}
\label{sec:esp_busqueda}

La mayor dificultad del Problema del Viajante reside en el enorme número de recorridos posibles que pueden generarse, incluso 
para instancias de tamaño medio. Este crecimiento del espacio de búsqueda es el principal motivo por el que la resolución
exacta del problema resulta inviable cuando aumenta el número de ciudades.

Supongamos una instancia formada por $n$ ciudades. Si se selecciona arbitrariamente una de ellas como ciudad de origen, 
$n-1$ ciudades pueden visitarse en cualquier orden, dando lugar a:

\begin{equation}
  \label{eq:eq5}
    (n-1)!
\end{equation}

posibles recorridos diferentes. Sin embargo, en el caso del Problema del Viajante simétrico, recorrer un ciclo en sentido 
contrario produce exactamente la misma solución, ya que el coste de desplazamiento entre dos ciudades es idéntico en ambos
sentidos. Por tanto, cada recorrido se contabiliza dos veces, obteniéndose finalmente:

\begin{equation}
  \label{eq:eq6}
    N=\frac{(n-1)!}{2}
\end{equation}

La tabla~\ref{tab:esp_busq} indica como evoluciona el número de posibles soluciones conforme aumenta el número de
ciudades. Se puede observar que el crecimiento factorial del espacio de búsqueda hace que incluso instancias relativamente
pequeñas sean imposibles de explorar exhaustivamente mediante fuerza bruta.

\begin{table}[ht]
    \renewcommand{\arraystretch}{1.3}
    \caption{Número de recorridos posibles para instancias simétricas.}
    \label{tab:esp_busq}
    \begin{center}
        \begin{tabular}{|c|c|}
            \hline
            \textbf{Número de ciudades} & \textbf{Recorridos posibles} \\
            \hline
            5 & 12 \\
            \hline
            10 & 181\,440 \\
            \hline
            15 & 43\,589\,145\,600 \\
            \hline
            20 & $6.08 \times 10^{16}$ \\
            \hline
            25 & $3.10 \times 10^{23}$ \\
            \hline
        \end{tabular}
    \end{center}
\end{table}

La figura~\ref{fig:esp_busq} muestra una comparativa de este crecimiento factorial del TSP frente a la complejidad
cuadrática y la exponencial, para reflejar de un modo más visual la diferencia abismal entre ellas.

\begin{figure}[H]
  \centering 
  {\includegraphics[width=0.85\linewidth]{tsp_factorial_growth.png}}
  \caption{Crecimiento factorial del TSP.}
  \label{fig:esp_busq}
\end{figure}

Este crecimiento factorial explica por qué los algoritmos exactos únicamente resultan prácticos para instancias de tamaño 
muy reducido. A medida que aumenta el número de ciudades, el tiempo necesario para explorar todas las soluciones posibles se 
vuelve inasequible, incluso para equipos de altas prestaciones. Como consecuencia, una parte de la investigación en optimización
combinatoria se ha orientado hacia el desarrollo de algoritmos aproximados, capaces de obtener soluciones de alta calidad sin 
necesidad de recorrer todo el espacio de búsqueda. Entre estas técnicas destacan las metaheurísticas, cuyo funcionamiento se 
va a desarrollar en los capítulos siguientes.

\section{Representación de una solución}
\label{sec:soluciones}

Una vez formulado el problema, es necesario definir cómo se representa una solución candidata. Esta representación constituye
la estructura sobre la que operan la mayoría de algoritmos de optimización, con especial interés en las heurísticas y 
metaheurísticas desarrolladas en los capítulos posteriores. 

En el Problema del Viajante, una solución puede representarse mediante una permutación de los vértices del grafo. Cada elemento
de la secuencia indica el orden en el que se visita una ciudad, asumiendo que, una vez alcanzado el último vértice, el recorrido
vuelva automáticamente al punto de partida:

\begin{equation}
  \label{eq:eq7}
    \pi = (1,4,2,5,3)
\end{equation}

O también utilizando flechas:

\begin{equation}
  \label{eq:eq8}
    1 \xrightarrow{} 4 \xrightarrow{} 2 \xrightarrow{} 5 \xrightarrow{} 3 \xrightarrow{} 1
\end{equation}

\begin{figure}[H]
  \centering 
  {\includegraphics[width=0.4\linewidth]{tsp_perm.png}}
  \caption{Solución candidata de una instancia pequeña.}
  \label{fig:tsp_perm}
\end{figure}

Consecuentemente, dos permutaciones distintas pueden representar el mismo recorrido. Por ejemplo:

\begin{equation}
  \label{eq:eq9}
    \pi = (4,2,5,3,1)
\end{equation}

Y la permutación~\ref{eq:eq8} describen exactamente el mismo ciclo, ya que únicamente cambia el vértice desde el que comienza
el recorrido.

A diferencia de otros problemas de optimización, cualquier solución válida debe contener exactamente una aparición de cada 
ciudad. Por tanto, las operaciones realizadas por los algoritmos de optimización no pueden modificar arbitrariamente la 
secuencia, sino que deben preservar siempre esta propiedad.

Otro factor importante es la calidad de la solución, la cual viene determinada por su función objetivo definida en la
ecuación~\ref{eq:eq2} como la suma de los costes asociados a todas las aristas recorridas.

Sobre esta representación se construyen la mayoría de operadores utilizados en optimización combinatoria. Intercambiar dos
ciudades, invertir un segmento del recorrido o combinar dos soluciones distintas son ejemplos de operaciones fundamentales
que se utilizan en muchas heurísticas y metaheurísticas para explorar el espacio de búsqueda.

\section{Aplicaciones del problema}
\label{sec:aplicaciones}

Aunque el Problema del Viajante nació como un problema puramente matemático, en la actualidad es uno de los modelos de
optimización combinatoria con mayor número de aplicaciones prácticas. Muchas situaciones reales pueden formularse como un TSP o
alguna de sus variantes, siempre que exista la necesidad de determinar el orden óptimo en el que deben visitarse un conjunto de
ubicaciones y minimizando un determinado coste.

Una de sus aplicaciones más conocidas se encuentra en el entorno de la logística y el transporte. Problemas como la distribución 
de mercancías, la planificación de autobuses escolares, la recogida de depósitos bancarios o la asignación de múltiples equipos
de trabajo pueden modelarse mediante variantes del problema, especialmente el MTSP. En estos escenarios, el objetivo 
consiste en minimizar, de forma paralela, la distancia recorrida y los costes operativos asociados a cada recorrido.

Otra aplicación aparece en los procesos industriales. En la fabricación de \ac{PCB}, por ejemplo, la cabeza de perforación debe
recorrer cientos o miles de posiciones realizando el menor número posible de desplazamientos improductivos. Del mismo modo, la 
generación de máscaras fotográficas para circuitos integrados, o el control de máquinas de corte y grabado pueden establecerse 
como problemas equivalentes al TSP, donde optimizar el recorrido implica una reducción directa del tiempo de fabricación.

Las aplicaciones del problema también se extienden a campos como la robótica, la topografía mediante \ac{GNSS}, la inspección
automática y las telecomunicaciones. En todos estos casos, resulta necesario planificar el recorrido óptimo de uno o varios 
agentes que deben visitar un conjunto de ubicaciones, manteniendo así un equilibrio entre la calidad de la solución y el tiempo
de cálculo requerido \cite{matai2010}. 

\begin{figure}[H]
  \centering 
  \fbox{
\includegraphics[width=0.7\linewidth]{tsp_applications.png}}
  \caption{Algunas aplicaciones prácticas del TSP.}
  \label{fig:applications}
\end{figure}

Esta amplia diversidad de aplicaciones explica el interés continuo que el Problema del Viajante sigue despertando en la comunidad
científica, así como el constante desarrollo de nuevas técnicas heurísticas y metaheurísticas capaces de resolver instancias 
cada vez mayores.



%%% Local Variables:
%%% TeX-master: "../book"
%%% End:
